\section{TFP Level and TFP Uncertainty}

Let \(a_t=\log A_t\). The level innovation to TFP has stochastic volatility:
\begin{align}
    a_{t+1}
    &=
    \rho_a a_t
    +
    \exp(V_{A,t})\sigma_A\varepsilon_{A,t+1},
    \qquad
    \varepsilon_{A,t+1}\sim N(0,1),
    \label{eq:tfp_level}
    \\
    V_{A,t+1}
    &=
    (1-\rho_{VA})\bar{V}_A
    +
    \rho_{VA} V_{A,t}
    +
    \sigma_{VA}\varepsilon_{VA,t+1},
    \qquad
    \varepsilon_{VA,t+1}\sim N(0,1).
    \label{eq:tfp_vol}
\end{align}
The primitive first-moment shock is \(\varepsilon_{A,t+1}\). The primitive
uncertainty shock is \(\varepsilon_{VA,t+1}\), which moves the conditional
volatility state.

Conditional on date \(t\) information,
\begin{equation}
    \Var_t(a_{t+1})=\sigma_A^2\exp(2V_{A,t}).
    \label{eq:conditional_tfp_variance}
\end{equation}
Define the centered variance state
\begin{equation}
    q_t
    \equiv
    \sigma_A^2\left[\exp(2V_{A,t})-\exp(2\bar{V}_A)\right].
    \label{eq:variance_state}
\end{equation}

The key conceptual point is that \(A_t\) is already the efficiency wedge. Hence
\(V_{A,t}\) is not a fifth wedge. It is a time-varying second moment of the
efficiency wedge:
\[
    \varepsilon_{VA,t}
    \quad \Rightarrow \quad
    \text{higher conditional variance of } A_{t+1}.
\]
The question is whether this second-moment movement in the efficiency wedge
creates risk correction terms that, when represented in a first-order BCA
accounting system, look like labor, investment, or resource wedges.
